% =========================
% CHAPTER 4 — METHODOLOGY + PRELIMINARY WORK + WORKPLAN
% File: chapter45-listing.tex
% =========================

\chapter{Methodology, preliminary work and workplan}\label{chap:ch4}

This chapter summarises the preparatory work completed to date and presents a concise, execution-oriented plan for the dissertation phase.
The core practical motivation is that, in heterogeneous coastal environments, na\"ively assimilating measurements into a single global
model may improve estimates locally while degrading predictions elsewhere, motivating \emph{regime-aware} inference and update logic
\parencite{bernacchi2025closingloop,fossum2020compact,ge2023watermasses}.

% ==========================================================
\section{Preliminary work (summary)} \label{sec:prelim_work}

\subsection{State of the art consolidation}
A structured synthesis of the literature was conducted to map how closed-loop adaptive sampling systems connect sensing, inference,
planning, multi-agent coordination, and model update. The main output is the taxonomy in Chapter~\ref{chap:ch3}, anchored in adaptive
sampling surveys \parencite{hwang2019auvreview} and representative works on closed-loop forecast improvement
\parencite{bernacchi2025closingloop}, compact representations for regime-like structure discovery \parencite{fossum2020compact},
profile-driven water-mass classification \parencite{ge2023watermasses}, flow-aware planning \parencite{aguiar2018trajectoryopt},
and multi-agent learning for plume monitoring \parencite{dalfabbro2025douroplume}. Classic work on coordinated sampling was used to
ground operational constraints \parencite{das2012coordinated,curtin2009aosn}.

\subsection{Problem framing and pipeline interface}
Following the review meeting, ``closing the loop'' is treated as an \emph{interface problem}: the dissertation does not aim to develop
a full operational data-assimilation system. Instead, the focus is on (i) inferring \emph{regimes} from in situ data, (ii) using regime
beliefs to guide multi-agent decisions, and (iii) evaluating whether regime-conditioned updates reduce harmful cross-regime coupling
\parencite{bernacchi2025closingloop,fossum2020compact,ge2023watermasses}.

\subsection{Scenarios and evaluation ladder}
Candidate scenarios were selected to support a validation ladder from controlled tests (enabling clean ablations and quantitative metrics)
to heterogeneous coastal cases (fronts/plumes) stressing regime transitions and robustness. The Douro plume setting is a natural reference
for multi-day dynamics and coordination constraints \parencite{dalfabbro2025douroplume}. Regime inference is centred on physically grounded
T/S/depth structure \parencite{ge2023watermasses}, with compact models as an unsupervised alternative \parencite{fossum2020compact}.

% ==========================================================
\section{Methodology overview and experimental design} \label{sec:method_overview}

\subsection{End-to-end loop (scope-aware)}
The dissertation methodology is organised as a lightweight closed-loop that remains compatible with operational constraints and the project scope.
At each coordination cycle, the system: (i) aggregates available measurements, (ii) infers a \emph{regime belief} over the domain (or over the
vehicle trajectory), (iii) scores candidate sampling actions using information/uncertainty proxies, and (iv) produces either a coordinated
allocation (when communication allows) or local decisions otherwise. The ``update'' component is intentionally implemented as a principled
\emph{proxy} (e.g., a belief refinement or a posterior map update) rather than full data assimilation, enabling controlled comparisons while
preserving the interface assumption emphasised in the review meeting \parencite{bernacchi2025closingloop,duarte2025geostatistical}.

\subsection{Regime inference: baselines and alternatives}
Regime inference is treated as a probabilistic classification or clustering problem over in situ observations.
The baseline will be a \emph{GMM/PCM-like} approach over profile-derived features (T/S/depth, stratification proxies, gradients), chosen for
its interpretability and low computational cost \parencite{ge2023watermasses}. Outputs are \emph{belief distributions} (rather than hard labels),
supporting boundary awareness (high uncertainty near transitions) and providing a natural signal to trigger regime-aware behaviour.

As an alternative representation, compact codes will be evaluated for unsupervised regime discovery and separability, inspired by dictionary learning
and sparse coding \parencite{fossum2020compact}. Because the dissertation targets profile data rather than synoptic surface imagery, the methodology
explicitly tests input designs (e.g., windowed profiles, derivative features, normalisation) and validates whether compact embeddings preserve physically
meaningful regime structure under distribution shift.

\subsection{Flow context: feasibility and boundary hypotheses from currents}
When velocity information is available (from numerical forecasts or onboard estimates), the methodology incorporates currents in two ways:
(i) \emph{feasibility and cost modelling} (time/energy-aware travel under time-varying flows) following flow-aware trajectory optimisation principles
\parencite{aguiar2018trajectoryopt}; and (ii) \emph{contextual regime cues} where flow-derived structures provide additional hypotheses for boundaries
(e.g., separation between coherent structures and strain-dominated regions). This complements T/S-driven regimes and supports tests of hybrid boundary
logic, where CTD-driven belief is combined with flow-derived boundary likelihood to prioritise sampling near moving interfaces under sparse measurements.

\subsection{Regime-aware sampling and update logic}
Regime beliefs are used to influence decisions in three concrete, testable ways:
\begin{enumerate}
  \item \textbf{Boundary emphasis:} increase utility near locations with high regime uncertainty (ambiguous class membership), treating transitions as
  first-class targets rather than incidental effects.
  \item \textbf{Team allocation across regimes:} encourage multi-AUV coverage diversity by allocating effort across regimes to reduce blind spots and
  minimise redundant sampling within a single regime under intermittent coordination \parencite{hwang2019auvreview,dalfabbro2025douroplume}.
  \item \textbf{Update gating/weighting:} condition the update proxy on regime belief to reduce harmful cross-regime coupling (e.g., downweight observations
  whose regime is inconsistent with the model component being updated, or maintain separate regime-conditioned local surrogates).
\end{enumerate}
These choices are designed to be evaluated via ablations (regime-agnostic vs regime-aware), without requiring a full DA pipeline while still testing the
central hypothesis motivating the dissertation.

% --------------------------
\subsection{Evaluation protocol, metrics, and ablations} \label{sec:metrics_ablations}
To ensure comparability, experiments will follow a pre-defined evaluation protocol that reports both \emph{within-regime} and \emph{cross-regime} behaviour.
In particular, the methodology distinguishes: (i) how well regimes are inferred (separability/calibration), (ii) how efficiently the team samples (coverage/redundancy),
and (iii) whether update proxies improve predictive consistency without inducing degradation elsewhere.

\begin{table}[H]
  \centering
  \caption{Planned evaluation axes, core metrics, and primary ablations.}
  \label{tab:eval_axes}
  \scriptsize
  \renewcommand{\arraystretch}{1.15}
  \setlength{\tabcolsep}{4pt}
  \begin{tabularx}{\linewidth}{>{\RaggedRight\arraybackslash}p{0.22\linewidth} X >{\RaggedRight\arraybackslash}p{0.25\linewidth}}
    \hline
    \textbf{Axis} & \textbf{What is measured} & \textbf{Ablations (examples)} \\
    \hline
    Regime inference &
    Separability and uncertainty behaviour (e.g., boundary regions show higher entropy); sensitivity to shift across scenarios &
    GMM/PCM-like vs compact codes; feature sets; calibration on/off \\
    \hline
    Informative sampling &
    Utility achieved per travel cost; boundary targeting vs interior redundancy; time-to-reach informative zones &
    Regime-agnostic utility vs boundary-weighted utility; myopic vs horizon-limited scoring \\
    \hline
    Multi-agent coordination &
    Redundancy/overlap under intermittent comms; robustness to dropouts; allocation stability across coordination windows &
    Centralised PCVRP-like baseline vs decentralised allocation; dropout/latency stress tests \\
    \hline
    Update behaviour (proxy) &
    Improvement within regime vs degradation elsewhere; consistency across regimes; sensitivity to mixed-regime trajectories &
    Regime-agnostic update vs regime-gated/weighted update; single-model vs regime-conditioned local surrogates \\
    \hline
    Flow awareness &
    Feasibility, arrival-time error, and cost under time-varying currents; impact on rendezvous/coordination windows &
    Flow-agnostic cost vs flow-aware cost; forecast mismatch margins \\
    \hline
  \end{tabularx}
\end{table}

\noindent The evaluation design explicitly addresses the review meeting emphasis on \emph{comparisons with ablations}: for each scenario, the dissertation will report
paired results for regime-agnostic and regime-aware variants under matched communication and flow assumptions, using consistent metrics and identical budgets.

% ==========================================================
\section{Workplan (Gantt) and risk management} \label{sec:workplan}

In line with the review meeting, the dissertation plan is expressed as a small set of parallel threads, with details tracked through deliverables rather than long narrative text.
Figure~\ref{fig:gantt_compact} summarises the workplan.

% IMPORTANT:
% Move \newcolumntype definitions to the preamble (ONLY ONCE), e.g.:
% \usepackage{array,tabularx,ragged2e}
% \newcolumntype{C}[1]{>{\centering\arraybackslash}p{#1}}
% \newcolumntype{L}[1]{>{\RaggedRight\arraybackslash}p{#1}}

\begin{figure}[H]
\centering

\makebox[\textwidth][c]{%
  \hspace*{-5cm}% <--- ajusta aqui: -0.3cm, -0.5cm, -0.7cm...
  \resizebox{\textwidth}{!}{%
    \begin{ganttchart}[
        time slot unit=month,
        time slot format=isodate-yearmonth,
        vgrid,
        hgrid,
        x unit=1.05cm,
        y unit chart=0.52cm,
        bar height=0.45,
        title height=0.75,
        title label font=\scriptsize,
        bar label font=\scriptsize,
        bar label node/.append style={align=right, text width=6.0cm},
        bar/.append style={fill=black!15, draw=black},
        milestone/.append style={fill=black!35, draw=black},
      ]{2026-02}{2026-06}
      \gantttitlecalendar{year, month=shortname} \\

      \ganttbar{Review state of the art}{2026-02}{2026-03} \\
      \ganttbar{Domain research}{2026-02}{2026-03} \\
      \ganttbar{ML on collected data}{2026-03}{2026-06} \\
      \ganttbar{Comparison of methods}{2026-05}{2026-06} \\
      \ganttbar{Writing}{2026-02}{2026-06} \\
      \ganttbar{Risk analysis \& mitigation}{2026-02}{2026-06} \\
      \ganttmilestone{Submission freeze}{2026-06} \\
    \end{ganttchart}%
  }%
}

\caption{Compact workplan (Gantt). Tasks are shown as parallel threads, with risk management running throughout.}
\label{fig:gantt_compact}
\end{figure}


\subsection{Risk matrix (actionable)}
Table~\ref{tab:risk_matrix} provides a compact risk matrix used to monitor the project throughout execution. Likelihood and impact are rated on a five-level qualitative scale
(from \emph{Rare} to \emph{Almost certain}, and from \emph{Very Low} to \emph{Very High}). Each risk is tracked with mitigation actions and periodic review checkpoints.
In particular, risks R1--R3 are treated as priority risks because they directly affect scientific validity (scenario realism and regime heterogeneity),
model robustness under shift, and operational assumptions (communication constraints).


\begin{table}[H]
  \centering
  \caption{Risk register (compact). Levels match Table~\ref{tab:risk_matrix}.}
  \label{tab:risk_register}
  \scriptsize
  \renewcommand{\arraystretch}{1.12}
  \setlength{\tabcolsep}{4pt}
  \begin{tabularx}{\linewidth}{C{0.07\linewidth} X C{0.12\linewidth} C{0.12\linewidth} L{0.30\linewidth}}
    \hline
    \textbf{ID} & \textbf{Risk} & \textbf{Lik.} & \textbf{Impact} & \textbf{Mitigation} \\
    \hline
    R1 & Low scenario realism / weak heterogeneity & Likely   & Very High & Freeze cases early; add controlled testbed. \\
    R2 & Regime model fails under shift            & Likely   & High      & Probabilistic + calibration; keep GMM/PCM baseline. \\
    R3 & Comms assumptions unrealistic             & Likely   & High      & Dropout/latency tests; hybrid coordination; track overlap. \\
    R4 & Unclear comparisons (metrics/ablations)   & Possible & Medium    & Predefine metrics; in- vs out-of-regime reporting. \\
    R5 & Update/loop closure exceeds scope         & Possible & Medium    & Use lightweight proxies; DA as external interface. \\
    R6 & Schedule slippage (impl. overhead)        & Possible & Low       & Continuous writing; monthly checkpoints; scope freeze. \\
    \hline
  \end{tabularx}
\end{table}

\begin{table}[t]
  \centering
  \caption{Risk matrix (Likelihood vs Impact). Cells are colour-coded (green = low, red = critical) and contain risk IDs (see Table~\ref{tab:risk_register}).}
  \label{tab:risk_matrix}
  \scriptsize
  \setlength{\tabcolsep}{6pt}
  \renewcommand{\arraystretch}{1.45}
  \begin{tabular}{l|C{2.0cm}C{2.0cm}C{2.0cm}C{2.0cm}C{2.0cm}}
    \hline
    \multicolumn{1}{c|}{} & \multicolumn{5}{c}{\textbf{Impact}} \\
    \textbf{Likelihood} & \textbf{Very Low} & \textbf{Low} & \textbf{Medium} & \textbf{High} & \textbf{Very High} \\
    \hline
    \textbf{Almost certain} &
    \cellcolor{orange!35} & \cellcolor{orange!45} & \cellcolor{red!35} &
    \cellcolor{red!55} & \cellcolor{red!75} \\
    \textbf{Likely} &
    \cellcolor{yellow!35} & \cellcolor{orange!30} & \cellcolor{orange!45} &
    \cellcolor{red!35}\textbf{R2, R3} & \cellcolor{red!55}\textbf{R1} \\
    \textbf{Possible} &
    \cellcolor{green!25} & \cellcolor{yellow!30}\textbf{R6} &
    \cellcolor{orange!30}\textbf{R4, R5} & \cellcolor{orange!45} & \cellcolor{red!35} \\
    \textbf{Unlikely} &
    \cellcolor{green!20} & \cellcolor{green!25} & \cellcolor{yellow!30} &
    \cellcolor{orange!30} & \cellcolor{orange!45} \\
    \textbf{Rare} &
    \cellcolor{green!20} & \cellcolor{green!20} & \cellcolor{green!25} &
    \cellcolor{yellow!30} & \cellcolor{orange!30} \\
    \hline
  \end{tabular}
\end{table}

% ==========================================================
\section{Summary} \label{sec:ch45_summary}

This chapter summarised the preparatory work and presented a scope-aware methodology and workplan aligned with the review meeting.
The dissertation will prioritise ML-based regime inference from in situ data and will evaluate its impact on multi-agent sampling decisions
and on regime-conditioned update behaviour in heterogeneous environments \parencite{fossum2020compact,ge2023watermasses,bernacchi2025closingloop}.
A concise Gantt chart captures the main execution threads, while an actionable risk matrix and an evaluation/ablation plan are provided to
support systematic comparisons under realistic communication and flow constraints.





% =========================
% FINAL CONSIDERATIONS / CONCLUSIONS
% (place in your conclusions chapter/section file)
% =========================

\chapter{Final Considerations and Conclusions} \label{chap:conclusions}

This report framed the dissertation topic of \emph{Dynamic Sampling with Multi-Agent AUVs Using Machine Learning} and reviewed the state of the art to support a research plan that is both scientifically grounded and operationally realistic. Overall, the literature converges on closed-loop adaptive sampling as the dominant paradigm, but it also shows that \emph{forecast improvement} is harder than pure mapping because it requires a consistent loop from uncertainty estimates to planning and back to model update. A central challenge in realistic coastal deployments is \emph{heterogeneity}: regime transitions (water masses, fronts, plume-influenced waters) can break naïve update assumptions, motivating regime-aware inference and regime-conditioned update strategies. The review also reinforces that uncertainty (or robust proxies) is critical to drive informative objectives and to reduce redundancy in multi-AUV missions under intermittent communications.

Based on these findings, the dissertation will prioritise \emph{probabilistic regime inference} from T/S/depth measurements (GMM/PCM-like) complemented by one alternative representation (e.g., compact codes/sparse coding) to benchmark robustness and separability. Planning will rely on informative objectives linked to uncertainty/information proxies and will explicitly test regime-aware criteria that emphasise ambiguous boundaries. For multi-agent operation, the approach will combine a decentralised allocation mechanism suitable for sparse communications with a strong centralised VRP-like baseline when coordination opportunities exist. The update component will remain scope-aware by using lightweight, principled update proxies, enabling systematic comparisons between regime-agnostic and regime-conditioned updates in controlled and coastal stress-test scenarios.

Key constraints and risks remain related to data representativeness and limited labels, distribution shift across seasons/locations, communication latency/dropouts, and onboard compute limits. These will be addressed through early scenario freezing, explicit robustness tests, communication stress testing, and continuous runtime/memory profiling. Immediate next steps are to lock the evaluation scenarios, implement end-to-end baselines with flow-aware costs, prototype and validate regime inference offline before integrating online belief updates, add a coordination baseline and quantify redundancy, and then run controlled experiments comparing regime-agnostic versus regime-conditioned update behaviour using pre-defined metrics and ablations.

